\documentclass[10pt,a4paper]{amsart}
\usepackage[margin=19mm]{geometry}
\usepackage{amsmath,amssymb,mathtools}
\usepackage[scr=rsfs,cal=euler]{mathalfa}
\usepackage{xfrac}
\usepackage[unicode,hidelinks,bookmarks=false]{hyperref}
\newcommand{\ie}{i.e.\ }
\newcommand{\cf}{cf.\ }
\newcommand{\resp}{respectively\ }
\newcommand{\Subsection}{\subsection}
\providecommand{\nobreakdash}{\nobreak\hbox{-}}
\setlength{\emergencystretch}{2em}
\multlinegap0pt
\providecommand{\R}{\mathbb R}
\usepackage{amssymb}
\usepackage{tikz}
\usepackage[shortlabels]{enumitem}
\newtheorem{lemma}{Lemma}
\newtheorem{corollary}{Corollary}
\providecommand{\wh}{\widehat}
\title{A continuum of Small-cap decouplings and Exponential Sums for the Moment Curve in $\R^4$}
\author{Jacob Glidewell}
\date{Source: arXiv:2605.27065v1 (CC BY 4.0)}
\begin{document}
\maketitle

\noindent\emph{Source and licence.} A continuum of Small-cap decouplings and Exponential Sums for the Moment Curve in $\R^4$. Version arXiv:2605.27065v1 (CC BY 4.0), distributed by arXiv under the Creative Commons Attribution 4.0 International licence, http://creativecommons.org/licenses/by/4.0/, as recorded in the arXiv licence metadata for this exact version and shown on the abstract page https://arxiv.org/abs/2605.27065v1. Full licence notice: \texttt{licenses/arxiv-2605.27065-CC-BY-4.0.txt}.\par\smallskip

\noindent\textit{Editorial scope.} Counted unit: the article's corollary cor:Lambda\_kEstimate (source0.tex lines 418--421). The article's complete original proof of it is delivered with it (source0.tex lines 423--437, one \texttt{proof} environment of 15 lines). No mathematics is written, repaired or re-derived: the statement and the proof are the article's own lines, in the article's own notation, and no retained line is substituted or re-flowed.  Retained uncounted as the source-local context the proof needs: lines 261--271; lines 281--285.  Nothing was omitted from any retained span, so the package carries no omission record.  No retained line contains a citation, so the wrapper carries no bibliography and no bibliography entry.  No retained line contains a figure environment, an image-inclusion command or drawing code, so no image is delivered and the package declares no asset dependency.  Editorial formatting only, in addition: an AMS \texttt{amsart} preamble that restates the article's own declarations and macros used by the retained lines, namely the environments \texttt{lemma}, \texttt{corollary} on separate counters, together with the standard packages those lines need.  The retained environments print in this document as Lemma 1, Corollary 1 in order of appearance, whereas the article prints the same items with the numbering of its own full document.  The complete native source of the article as distributed in the arXiv e-print for this exact version is bundled unchanged as \texttt{sources/201467/source0.tex}.
\begin{lemma}\label{lem:PruningProperties}
For each $\gamma_k$,
    \begin{enumerate}
        \item[a)] 
        $$|F^{(k)}_{\gamma_k}(x)|\lesssim |F^{(k+1)}_{\gamma_k}(x)|\le \sum_{\gamma_{k+1}\subset \gamma_k}|F_{\gamma_{k+1}}^{(k+1)}(x)|,$$

        \item[b)]  $$\|F_{\gamma_k}^{(k)}\|_\infty \lesssim \min \left\{\frac{R^{\beta}}{\lambda}, \frac{R^\beta}{R_k}\right\},$$

        \item[c)] $$\text{supp }\wh{F_{\gamma_k}^{(k+1)}}\subset \text{supp }\wh{F_{\gamma_k}^{(k)}}\subset C \gamma_k.$$
    \end{enumerate}
\end{lemma}

\begin{lemma}\label{lem:L6andL12estimates}
 For any scale $R_k\le R^{1/2}$, we have
    $$\|F_{\gamma_k}^{(k)}\|_{L^{6}(\Omega_R)}\lessapprox |\Omega_R|^{\frac{1}{6}}\left(\frac{R^{\beta}}{R_k}\right)^{1/2}.$$
 For $R_k\le R^{1/3}$, we have $$\|F_{\gamma_k}^{(k)}\|_{L^{12}(\Omega_R)}\lessapprox \min \left\{\frac{R^{\frac{3}{4}\beta}}{\lambda^{\frac{1}{2}}R_k^{\frac{1}{4}}}, \frac{R^{\frac{3}{4}\beta}}{R^{\frac{1}{12}}R_k^{\frac{1}{2}}}\right\}|\Omega_R|^{\frac{1}{12}}.$$ 
\end{lemma}

\begin{corollary}\label{cor:Lambda_kEstimate}
    For $4\le \frac{2}{\beta}\le q\le 6$, and $R_k=R^{2/q}\in [R^{1/3}, R^{\beta}]$, 
    $$|\Lambda_k|\lessapprox\min\left\{\frac{R^{\beta(q-3)+1}}{R_k^3\lambda^{2q-6}}, \frac{R^{\beta(q-3)+1}}{R_k^{2q-3}}\right\}|\Omega_R|.$$ For $R_k\le R^{1/3}$, $$|\Lambda_k|\lessapprox \min\left\{\frac{R^{3\beta}}{\lambda^6}, \frac{R^{3\beta}}{R_k^3R}\right\}|\Omega_R|.$$
\end{corollary}

\begin{proof}
For $R_k= R^{2/q}$ for $q\in [\frac{2}{\beta},6]$, then by Lemma \ref{lem:L6andL12estimates}\begin{align*}
    R^{\beta q}|\Lambda_k|\lesssim\int_{\Omega_R}|g_k^h|^q&\lessapprox R_k^{-1}R \sum_{\gamma_{k+1}}\int_{\Omega_R} |F_{\gamma_{k+1}}^{(k+1)}|^{2q}\\
    &\lesssim R_k^{-1}R\left(\frac{R^\beta}{\lambda}\right)^{2q-6}\sum_{\gamma_{k+1}}\int_{\Omega_R}|F_{\gamma_{k+1}}^{(k+1)}|^{6}\\
    &\lessapprox R_k^{-1}R \left(\frac{R^\beta}{\lambda}\right)^{2q-6} R_k\left(\frac{R^{\beta}}{R_k}\right)^3|\Omega_R|\\
    &= \frac{R^{\beta(2q-3)+1}}{R_k^3\lambda^{2q-6}}|\Omega_R|.
\end{align*} The second bound in this case is found by repeating the same computation but using Lemma \ref{lem:PruningProperties} $$\|F_{\gamma_{k+1}}^{(k+1)}\|_\infty \lesssim \frac{R^{\beta}}{R_k}.$$
%
For $R_k\le R^{1/3}$, \begin{align*}
    R^{6\beta}|\Lambda_k|\lesssim\int_{\Omega_R}|g_k^h|^6&\lessapprox \left(\sum_{\gamma_{k+1}}\|F_{\gamma_{k+1}}^{(k+1)}\|_{L^{12}(\Omega_R)}^4\right)^3\\
    &\lessapprox  \left(\sum_{\gamma_{k+1}}\min \left\{\frac{R^{3\beta}}{\lambda^2R_k}, \frac{R^{3\beta}}{R^{\frac{1}{3}}R_k^2}\right\}|\Omega_R|^{\frac{1}{3}}\right)^3\\
    &\lessapprox \min \left\{\frac{R^{9\beta}}{\lambda^6}, \frac{R^{9\beta}}{RR_k^3}\right\}|\Omega_R|
\end{align*}

\end{proof}

\end{document}
